%!TEX root = ../thesis.tex
%*******************************************************************************
%*********************************** Chương 1 **********************************
%*******************************************************************************

\chapter{GIỚI THIỆU}
\label{chap:intro}

\ifpdf
    \graphicspath{{Chapter1/Figs/Raster/}{Chapter1/Figs/PDF/}{Chapter1/Figs/}}
\else
    \graphicspath{{Chapter1/Figs/Vector/}{Chapter1/Figs/}}
\fi

% ------------------------------ Từ viết tắt -----------------------------------
% Các mục dưới đây tự gom vào "Danh mục từ viết tắt" (nhóm Z).
\nomenclature[z-VQA]{VQA}{Visual Question Answering: hỏi đáp trên ảnh}
\nomenclature[z-VLM]{VLM}{Vision-Language Model: mô hình nhận cả ảnh và văn bản, sinh văn bản}
\nomenclature[z-OCR]{OCR}{Optical Character Recognition: nhận dạng chữ trong ảnh}
\nomenclature[z-PEFT]{PEFT}{Parameter-Efficient Fine-Tuning: chỉ huấn luyện một phần rất nhỏ các tham số}
\nomenclature[z-LoRA]{LoRA}{Low-Rank Adaptation: thêm hai ma trận hạng thấp vào các lớp tuyến tính và chỉ huấn luyện chúng}
\nomenclature[z-QLoRA]{QLoRA}{LoRA trên mô hình nền đã lượng tử hóa 4-bit (NF4)}
\nomenclature[z-NF4]{NF4}{4-bit NormalFloat: kiểu lượng tử hóa 4-bit}
\nomenclature[z-EM]{EM}{Exact Match: câu trả lời trùng khít với đáp án sau chuẩn hóa}
\nomenclature[z-VQA-Acc]{VQA-Acc}{VQA Accuracy: min(số đáp án gold trùng / 3, 1)}
\nomenclature[z-ANLS]{ANLS}{Average Normalized Levenshtein Similarity: độ giống theo ký tự}
\nomenclature[z-POPE]{POPE}{Polling-based Object Probing Evaluation: đo hallucination bằng câu hỏi có/không}
\nomenclature[z-SOTA]{SOTA}{State of the art: kết quả tốt nhất đã công bố}

%********************************** Mục 1.1 ************************************
\section{Bối cảnh và động lực}
\label{sec:intro-background}

\subsection{Hỏi đáp trên ảnh và vision-language model}

Visual Question Answering (VQA) yêu cầu hệ thống trả lời một câu hỏi bằng ngôn
ngữ tự nhiên dựa trên nội dung ảnh~\cite{antol2015vqa}. Khác với phân loại ảnh
hay phát hiện vật thể, VQA cần hiểu đồng thời ảnh, ngữ nghĩa câu hỏi và đôi khi
cả chữ xuất hiện trong ảnh~\cite{singh2019textvqa,biten2019stvqa}. Các hệ thống
VQA ban đầu ghép một bộ mã hóa ảnh (CNN) với một bộ mã hóa câu hỏi (LSTM) và
chọn đáp án từ một tập cố định. Các vision-language model (VLM) gần đây như
BLIP-2~\cite{li2023blip2}, LLaVA~\cite{liu2023llava},
Florence-2~\cite{xiao2024florence2} và Qwen2.5-VL~\cite{bai2025qwen25vl} thay
đổi cách tiếp cận: một bộ mã hóa ảnh biến ảnh thành chuỗi token, một mô hình
ngôn ngữ lớn đọc chuỗi đó cùng câu hỏi và sinh câu trả lời tự do. Các VLM này
trả lời được câu hỏi mở và đọc được chữ trong ảnh ở mức zero-shot.

\subsection{Tiếng Việt với các VLM thiên tiếng Anh}

Phần lớn VLM mã nguồn mở được tiền huấn luyện chủ yếu trên dữ liệu tiếng Anh và
tiếng Trung. Khi dùng cho tiếng Việt, nhóm quan sát bốn vấn đề riêng, ký hiệu
V1--V4 và dùng xuyên suốt khóa luận:

\begin{description}
  \item[V1, đọc chữ có dấu trong ảnh.] Dấu thanh là nét rất nhỏ; đọc sai một
  dấu là đổi nghĩa (``BÁNH MÌ'' đọc thành ``banh mi'').
  \item[V2, sai ngôn ngữ hoặc sai định dạng.] Mô hình trả lời bằng tiếng Anh,
  chen chữ Trung, hoặc trả lời cả câu dài trong khi đáp án chỉ vài từ.
  \item[V3, thứ tự chữ và ranh giới từ.] Khoảng trắng trong tiếng Việt tách âm
  tiết, không tách từ; đọc sai thứ tự các vùng chữ là vỡ từ (``cửa hàng điện
  thoại'' thành ``cửa điện hàng thoại'').
  \item[V4, hallucination.] Câu trả lời trôi chảy nhưng nói về vật hoặc chữ
  không có trong ảnh.
\end{description}

Với dữ liệu tiếng Việt, khoảng cách giữa máy và người còn lớn: trên
ViTextVQA~\cite{nguyen2024vitextvqa}, kết quả công bố tốt nhất trước khóa luận
này là EM 25,48, trong khi người đạt 88,00.

\subsection{Ràng buộc tài nguyên}

Huấn luyện lại toàn bộ một VLM vài tỉ tham số (full fine-tuning) cần nhiều GPU
dung lượng lớn, vượt quá điều kiện của một nhóm sinh viên. Trường không có máy
chủ GPU cho đề tài, nên mọi thí nghiệm chạy trên Google Colab. Đề tài vì vậy đặt
hiệu quả thành một ràng buộc thiết kế tường minh: chỉ dùng parameter-efficient
fine-tuning (LoRA~\cite{hu2022lora}, QLoRA~\cite{dettmers2023qlora}) và báo cáo
kèm thời gian huấn luyện, thời gian suy luận và bộ nhớ GPU cho mọi kết quả.

%********************************** Mục 1.2 ************************************
\section{Phát biểu bài toán}
\label{sec:intro-problem}

Cho ảnh $I$ và câu hỏi tiếng Việt $q$, hệ thống sinh câu trả lời tiếng Việt
$\hat{a} = f(I, q)$. Câu trả lời được so với đáp án gold $a$ sau cùng một bước
chuẩn hóa văn bản. Hai bộ dữ liệu đặt ra hai yêu cầu khác nhau:

\begin{itemize}
  \item \textbf{ViVQA}~\cite{tran2021vivqa}: câu hỏi chung về ảnh (màu sắc, số
  lượng, vật thể, vị trí); 99,83\% đáp án dài không quá 3 từ, nên EM và VQA-Acc
  là thước đo phù hợp.
  \item \textbf{ViTextVQA}~\cite{nguyen2024vitextvqa}: đáp án là chữ đọc được
  trong ảnh, trung bình 4,18 từ; ANLS~\cite{biten2019stvqa} được dùng thêm vì
  chịu được lỗi đọc nhỏ ở mức ký tự.
\end{itemize}

Với ví dụ của ViVQA, câu hỏi ``màu của xe đạp là gì?'' có đáp án ``màu đỏ''. Với
ví dụ của ViTextVQA, câu hỏi ``cửa hàng có biển màu vàng là nơi nào?'' có đáp án
``đại lý xổ số miền Bắc'', chữ lấy trực tiếp từ biển hiệu.

%********************************** Mục 1.3 ************************************
\section{Khoảng trống nghiên cứu}
\label{sec:intro-gaps}

Từ khảo sát tài liệu ở Chương~\ref{chap:related}, nhóm xác định bốn khoảng
trống:

\begin{description}
  \item[G1.] Tiền huấn luyện thiên lệch tiếng Anh: mô hình zero-shot trả lời bằng
  tiếng Anh và sai định dạng khi gặp câu hỏi tiếng Việt.
  \item[G2.] VQA tiếng Việt còn yếu, nhất là với scene-text: SOTA ViTextVQA là EM
  25,48 so với 88,00 của người.
  \item[G3.] Thích nghi VLM cho VQA tiếng Việt bằng PEFT chưa được nghiên cứu:
  các công trình VQA tiếng Việt trong nhóm khảo sát đều huấn luyện mô hình riêng
  hoặc dùng mô hình đóng băng, không dùng LoRA/QLoRA trên VLM lớn.
  \item[G4.] Chưa có cách đo hallucination cho VQA tiếng Việt: chưa có benchmark
  kiểu POPE~\cite{li2023pope} bằng tiếng Việt.
\end{description}

Khóa luận tập trung vào G3 và dùng G2, G4 làm nơi kiểm chứng.

%********************************** Mục 1.4 ************************************
\section{Mục tiêu}
\label{sec:intro-objectives}

Mục tiêu tổng quát là thích nghi một VLM mã nguồn mở cho VQA tiếng Việt dưới
ngân sách tham số nghiêm ngặt và kiểm chứng bằng thực nghiệm ít nhất một cải
tiến so với baseline đã tái lập. Mục tiêu này được chia thành các mục tiêu cụ
thể trong Bảng~\ref{tab:objectives}, mỗi mục tiêu có tiêu chí đo lường.

\begin{table}[htbp]
\caption{Mục tiêu và tiêu chí hoàn thành}
\label{tab:objectives}
\centering
\small
\begin{tabular}{p{0.30\textwidth}p{0.60\textwidth}}
\toprule
\textbf{Mục tiêu} & \textbf{Tiêu chí hoàn thành} \\
\midrule
Tái lập đúng một baseline &
Chạy lại Qwen2.5-VL-3B-Instruct trên một benchmark có số công bố và đạt sai lệch
nhỏ (TextVQA val: nhóm đạt 79,60 so với 79,3 công bố) \\
Đo baseline zero-shot trên dữ liệu tiếng Việt &
EM, VQA-Acc, ANLS trên toàn bộ test ViVQA và ViTextVQA, cùng một giao thức \\
Thích nghi chỉ bằng PEFT &
QLoRA 4-bit, số tham số huấn luyện $<1\%$, chạy trên một GPU Colab \\
Đề xuất và kiểm chứng $\geq 1$ cải tiến &
Cải tiến so với baseline dưới điều kiện y hệt, có kiểm định thống kê ghép cặp và
ablation \\
Phân tích lỗi và hallucination &
Phân rã lỗi theo loại câu hỏi và loại lỗi; đo hallucination kiểu POPE cho tiếng
Việt \\
Prototype web &
Tải ảnh, hỏi tiếng Việt, nhận câu trả lời kèm chữ OCR và mức tin cậy; cho cùng
kết quả với đường đánh giá \\
Tái lập được &
Mọi run ghi lại split, seed, phần cứng, siêu tham số, thời gian; gói tái lập gồm
code, adapter, config, script suy luận \\
\bottomrule
\end{tabular}
\end{table}

%********************************** Mục 1.5 ************************************
\section{Câu hỏi nghiên cứu và giả thuyết}
\label{sec:intro-rq}

\begin{description}
  \item[RQ1.] Zero-shot đạt bao nhiêu trên ViVQA và ViTextVQA, và lỗi tiếng Việt
  nào chiếm ưu thế?
  \item[RQ2.] QLoRA cải thiện bao nhiêu so với zero-shot, khi không dùng full
  fine-tuning?
  \item[RQ3.] Chữ OCR có dấu, đúng thứ tự đọc, có giúp mô hình đọc chữ trong ảnh
  không, và có cộng hưởng với QLoRA không?
  \item[RQ4.] Sau khi thích nghi bằng QLoRA, mô hình bịa nội dung nhiều hơn hay
  ít hơn?
\end{description}

Mỗi giả thuyết trong Bảng~\ref{tab:hypotheses} có ngưỡng định lượng được ghi lại
trước khi chạy thí nghiệm tương ứng.

\begin{table}[htbp]
\caption{Giả thuyết và câu hỏi nghiên cứu tương ứng}
\label{tab:hypotheses}
\centering
\small
\begin{tabular}{p{0.06\textwidth}p{0.64\textwidth}p{0.18\textwidth}}
\toprule
\textbf{ID} & \textbf{Giả thuyết} & \textbf{Trả lời} \\
\midrule
H1 & QLoRA trên Qwen2.5-VL-3B tăng EM ít nhất 5 điểm so với chính mô hình đó ở
chế độ zero-shot. & RQ2; V1--V3 \\
H2 & Chữ OCR có dấu, sắp xếp từ trái-trên đến phải-dưới, đưa vào prompt tăng
ANLS ít nhất 3 điểm trên ViTextVQA; kết hợp với QLoRA cho kết quả cao nhất. &
RQ3; V1, V3 \\
H3 & Prompt định dạng tiếng Việt (trả lời ngắn, có dấu) tăng EM ít nhất 3 điểm
và giảm câu trả lời sai ngôn ngữ hoặc mất dấu. & RQ1; V2 \\
H4 & QLoRA không làm tăng tỉ lệ hallucination so với zero-shot. & RQ4; V4 \\
\bottomrule
\end{tabular}
\end{table}

%********************************** Mục 1.6 ************************************
\section{Đóng góp}
\label{sec:intro-contrib}

Nhóm chọn Qwen2.5-VL-3B-Instruct~\cite{bai2025qwen25vl} làm baseline, tái lập
baseline này trước, rồi thích nghi nó cho tiếng Việt chỉ bằng QLoRA. Mô hình nền
3 tỉ tham số được lượng tử hóa 4-bit và giữ nguyên; nhóm chỉ huấn luyện 7,37
triệu tham số adapter, tức 0,196\% tổng số tham số, trên một GPU NVIDIA L4. Các
đóng góp của khóa luận gồm:

\begin{enumerate}
  \item \textbf{QLoRA nhận biết OCR.} Văn bản đọc được trong ảnh bằng VietOCR
  (giữ dấu, sắp theo thứ tự đọc) được đưa vào prompt ngay trong lúc huấn luyện,
  để adapter học cách dùng nó. Trên toàn bộ tập test ViTextVQA (10.028 câu),
  cách này đạt EM 53,35 so với 50,26 của QLoRA thông thường với cùng dữ liệu
  ($+3{,}09$ EM, kiểm định McNemar $p = 6{,}7\times10^{-21}$), và cao hơn kết quả
  công bố tốt nhất trên cùng tập test (ViTextBLIP-2, EM
  25,48~\cite{nguyen2024vitextvqa}).
  \item \textbf{Một cơ chế được tách bạch.} Thí nghiệm $2\times2$ cho thấy OCR
  chỉ có ích khi được dùng cả lúc huấn luyện lẫn lúc suy luận: đưa OCR vào prompt
  của một mô hình chưa từng thấy OCR không cải thiện kết quả.
  \item \textbf{Kết quả trên VQA tiếng Việt tổng quát.} Trên ViVQA, QLoRA nâng
  EM từ 22,09 (zero-shot) lên 76,64.
  \item \textbf{Bộ đo hallucination tiếng Việt.} POPE-vi, gồm 4.985 câu hỏi
  có/không trên 500 ảnh chưa thấy khi huấn luyện.
  \item \textbf{Phân tích và prototype đã kiểm chứng.} Ablation cho từng thành
  phần, phân tích lỗi tự động và duyệt tay, và một prototype web đã được kiểm
  chứng cho ra cùng câu trả lời với đường đánh giá chính thức.
\end{enumerate}

\subsection{Ý nghĩa của đề tài}

\paragraph{Ý nghĩa khoa học.} Khóa luận cung cấp một nghiên cứu có hệ thống về
PEFT cho VLM trên VQA tiếng Việt, gồm baseline tái lập, ablation cho từng thành
phần và kiểm định thống kê ghép cặp. Kết quả chính không chỉ là một con số cao
hơn mà còn là một cơ chế được tách bạch: lợi ích của OCR đến từ việc huấn luyện
mô hình dùng OCR, không đến từ việc thêm OCR vào prompt. Nhóm cũng đưa ra bộ
POPE tiếng Việt để đo hallucination và một phân tích có số đo cho từng vấn đề
V1--V4.

\paragraph{Ý nghĩa thực tiễn.} Hệ thống đọc và trả lời câu hỏi về biển hiệu,
nhãn hàng, bảng giá, thực đơn và biểu mẫu tiếng Việt. Đây là chức năng nền cho
công nghệ hỗ trợ người khiếm thị và cho các ứng dụng tra cứu thông tin từ ảnh. Vì
chỉ cần một GPU và một adapter 29,5~MB, cách làm có thể áp dụng cho các nhóm hoặc
tổ chức có tài nguyên hạn chế.

\paragraph{Ý nghĩa phương pháp.} Mọi kết luận trong khóa luận dựa trên luật
quyết định được ghi lại trước khi chạy, kiểm định ghép cặp trên cùng tập test,
và log đầy đủ (split, seed, GPU, thời gian). Quy trình này giúp tránh việc chọn
kết quả sau khi đã thấy số liệu, và cho phép người khác chạy lại từng con số từ
mã nguồn công khai.

%********************************** Mục 1.7 ************************************
\section{Phạm vi và giới hạn}
\label{sec:intro-scope}

\begin{table}[htbp]
\caption{Trong và ngoài phạm vi}
\label{tab:scope}
\centering
\small
\begin{tabular}{p{0.45\textwidth}p{0.45\textwidth}}
\toprule
\textbf{Trong phạm vi} & \textbf{Ngoài phạm vi} \\
\midrule
Hai bộ dữ liệu ViVQA và ViTextVQA, split cố định &
Xây bộ dữ liệu huấn luyện mới hoặc tăng cường dữ liệu \\
Một baseline Qwen2.5-VL-3B-Instruct; các VLM khác chỉ để tham khảo &
So sánh nhiều backbone như baseline chính \\
Thích nghi chỉ bằng PEFT (QLoRA 4-bit); chữ OCR làm đầu vào phụ &
Full fine-tuning; huấn luyện VLM từ đầu \\
Metric EM, VQA-Acc, ANLS (token-F1 làm chỉ số phụ); ablation, phân tích lỗi,
hallucination & Retrieval-augmented VQA với tri thức ngoài \\
Prototype web chạy trên Colab với một GPU &
Triển khai sản phẩm quy mô lớn, ứng dụng di động \\
\bottomrule
\end{tabular}
\end{table}

\paragraph{Giả định.} Nhóm giả định rằng split chính thức của hai bộ dữ liệu
đại diện cho phân bố câu hỏi thực tế của từng bộ, rằng đáp án gold là đúng ở đa
số câu (phần gold sai được đo và báo cáo riêng ở Chương~\ref{chap:results}), và
rằng GPU L4 trên Colab cho kết quả ổn định giữa các phiên khi cố định seed và
cấu hình.

\paragraph{Giới hạn biết trước.} Kết quả bị giới hạn bởi một mô hình nền 3 tỉ
tham số ở 4-bit và bởi chất lượng OCR; mọi thí nghiệm chính chạy với một seed,
riêng phần so sánh OCR được kiểm thêm trên ba seed. EM là thước đo khắt khe với
câu trả lời đúng ý nhưng khác cách viết, nên nhóm bổ sung ANLS và duyệt tay. Các
giới hạn này được phân tích lại ở Chương~\ref{chap:discussion} và
Chương~\ref{chap:conclusion}.

%********************************** Mục 1.8 ************************************
\section{Cấu trúc khóa luận}
\label{sec:intro-outline}

Chương~\ref{chap:related} tổng quan VQA, vision-language model, parameter-efficient
fine-tuning và các nghiên cứu VQA tiếng Việt. Chương~\ref{chap:methods} trình bày
dữ liệu, baseline, phương pháp QLoRA nhận biết OCR và giao thức đánh giá.
Chương~\ref{chap:results} báo cáo kết quả chính và các ablation.
Chương~\ref{chap:discussion} thảo luận phân tích lỗi, hallucination và giới hạn.
Chương~\ref{chap:conclusion} kết luận. Phụ lục~\ref{appendix:pm} trình bày kế
hoạch quản lý dự án.
